\section{Resultados e Discussão}

\subsection{Caracterização do grid search}

\begin{figure}[!ht]
  \centering
  \includegraphics[width=0.92\textwidth]{grid_search_summary_article.png}
  \caption{Erro médio das configurações no grid STO-3G. A linha pontilhada
  marca 1 kcal/mol. Bordô indica sucesso em todas as geometrias, azul em parte
  delas e cinza em nenhuma. Diferentes \textit{reps} foram agregadas.}
  \label{fig:grid-summary}
\end{figure}

A Figura~\ref{fig:grid-summary} mostra que o ranking depende do sistema. Em
H$_2$, combinações fisicamente inspiradas atingem erros muito inferiores ao
limiar em todas as geometrias. Em LiH, somente parte das combinações alcança
precisão química. Em BeH$_2$, mesmo as melhores possuem erro médio acima de
1 kcal/mol, embora resolvam algumas distâncias. Essa variação motiva um ranking
condicionado à instância, mas também antecipa a dificuldade de generalização.

\subsection{Predição do erro}

\begin{figure}[!ht]
  \centering
  \includegraphics[width=\textwidth]{predicted_vs_real_models.png}
  \caption{Erro previsto e real no protocolo agrupado por geometria. A linha
  tracejada representa predição ideal; ambos os eixos são logarítmicos.}
  \label{fig:predictions}
\end{figure}

\begin{table}[!ht]
\centering
\caption{Média dos cinco folds no protocolo agrupado por geometria.}
\label{tab:regression}
\begin{tabular}{lrrrr}
\toprule
Modelo & MAE & RMSE & MAE alvo log & $R^2$ alvo log \\
       & \multicolumn{2}{c}{(kcal/mol)} & & \\
\midrule
Dummy         & 37,980 & 74,302 & 1,571 & -0,015 \\
Ridge         & 24,335 & 53,039 & 0,809 &  0,677 \\
Random Forest & \textbf{17,316} & \textbf{35,799} & \textbf{0,580} & 0,765 \\
MLP           & 18,746 & 37,688 & 0,596 & \textbf{0,774} \\
\bottomrule
\end{tabular}
\end{table}

A Tabela~\ref{tab:regression} mostra que Random Forest obteve os menores MAE e
RMSE, enquanto a MLP apresentou o maior $R^2$ log. A
Figura~\ref{fig:predictions} indica que ambos capturam a separação entre erros
baixos e altos, mas comprimem extremos. A transformação logarítmica evita que
poucos erros de centenas de kcal/mol dominem o ajuste. O piso visual em
$10^{-6}$ decorre apenas do gráfico: predições negativas no alvo log foram
truncadas em zero e, para exibição em escala logarítmica, zero foi representado
por $10^{-6}$ kcal/mol.

\subsection{Qualidade das recomendações}

Somente 23 das 27 geometrias possuem ao menos um candidato dentro da precisão
química; 85,2\% é, portanto, o teto de cobertura deste grid.

\begin{table}[!ht]
\centering
\caption{Resultados top-$k$ agrupados por geometria. Erro e regret em
kcal/mol; a última coluna conta candidatos evitados na inferência.}
\label{tab:topk}
\begin{tabular}{llrrrr}
\toprule
Modelo & $k$ & Cobertura & Erro & Regret & Candidatos evitados \\
\midrule
Aleatório       & 1 & 26,9\% & 40,099 & 39,316 & 94,4\% \\
Melhor global   & 1 & 77,8\% &  1,410 &  0,627 & 94,4\% \\
Melhor/molécula & 1 & 74,1\% &  1,593 &  0,810 & 94,4\% \\
MLP             & 1 & 77,8\% &  1,433 &  0,650 & 94,4\% \\
\midrule
Aleatório       & 3 & 55,4\% &  7,074 &  6,291 & 83,3\% \\
Melhor global   & 3 & 77,8\% &  1,008 &  0,225 & 83,3\% \\
Melhor/molécula & 3 & 77,8\% &  1,008 &  0,225 & 83,3\% \\
Random Forest   & 3 & 77,8\% &  0,961 &  0,178 & 83,3\% \\
MLP             & 3 & \textbf{85,2\%} & \textbf{0,783} & \textbf{0,000} & 83,3\% \\
Oráculo         & 3 & 85,2\% & 0,783 & 0,000 & -- \\
\bottomrule
\end{tabular}
\end{table}

A Tabela~\ref{tab:topk} mostra que, no top-1, a MLP não supera o melhor
baseline global. No top-3, a maior cobertura observada foi da MLP: 85,2\%
contra 77,8\% dos baselines global e por molécula, diferença de apenas duas
geometrias. O bootstrap pareado da diferença foi 7,4 pontos percentuais, com
IC$_{95\%}$ de 0,0 a 18,5 por geometria e de 0,0 a 22,2 ao reamostrar moléculas
como clusters. Logo, os dados não demonstram vantagem estatística da MLP.

\begin{figure}[!ht]
  \centering
  \includegraphics[width=0.86\textwidth]{topk_success_vs_reduction.png}
  \caption{Cobertura top-$k$ versus fração de candidatos não executados na
  inferência. Os pontos da MLP identificam $k=1,3,5$; faixas representam
  intervalos bootstrap de 95\%.}
  \label{fig:topk}
\end{figure}

A Figura~\ref{fig:topk} apresenta o compromisso entre cobertura e quantidade de
candidatos. Para $k=3$, o recomendador executaria três em vez de 18 candidatos,
redução nominal de 83,3\%. Isso não equivale a 83,3\% de tempo ou custo
quântico: ansätze e otimizadores possuem profundidades e números de chamadas
diferentes. Além disso, produzir o modelo exigiu previamente as 486 execuções.
O benefício depende, portanto, de amortizar o treinamento em novas geometrias
relacionadas. O oráculo é apenas limite superior e requer conhecer o grid todo.

\subsection{Ablação e generalização}

\begin{table}[!ht]
\centering
\caption{Desempenho da MLP nos protocolos de generalização.}
\label{tab:ablation}
\begin{tabular}{lrrr}
\toprule
Protocolo & MAE alvo log & $R^2$ alvo log & Cobertura top-3 \\
\midrule
Geometrias, com molécula & 0,596 &  0,774 & 85,2\% \\
Geometrias, sem molécula & 0,629 &  0,754 & 77,8\% \\
Molécula não vista       & 1,274 & -0,104 & 77,8\% \\
\bottomrule
\end{tabular}
\end{table}

Como mostra a Tabela~\ref{tab:ablation}, remover o identificador molecular
reduziu a cobertura top-3 da MLP ao mesmo valor do baseline global. Ao deixar
uma molécula inteira fora do treino, o $R^2$ log tornou-se negativo. Assim,
parte relevante do sinal aprendido é um ranking típico das famílias moleculares
observadas, e não uma representação química geral.

O estudo usa somente três moléculas, STO-3G, uma semente, 100 iterações e
simulação ideal. O bootstrap por geometria subestima potencialmente a
dependência intramolecular; o intervalo por clusters usa apenas três unidades e
também é impreciso. O recomendador escolhe apenas candidatos presentes no grid
e não cria ansätze nem substitui a otimização dos parâmetros. Os resultados
demonstram viabilidade dentro do domínio estudado, não generalização para
hardware, bases maiores ou moléculas arbitrárias.
